\documentclass[10pt,a4paper]{letter}
\usepackage[top=1.8cm,bottom=1.8cm,left=2.4cm,right=2.4cm]{geometry}
\usepackage[colorlinks=true,urlcolor=blue,linkcolor=blue]{hyperref}

\signature{Quang-Vinh Dang\\British University Vietnam\\Hung Yen, Vietnam\\\href{mailto:vinh.dq4@buv.edu.vn}{vinh.dq4@buv.edu.vn}\\ORCID 0000-0002-3877-8024}
\address{Quang-Vinh Dang\\British University Vietnam\\Hung Yen, Vietnam\\\href{mailto:vinh.dq4@buv.edu.vn}{vinh.dq4@buv.edu.vn}}
\date{13 November 2026}

\begin{document}

\begin{letter}{The Editors-in-Chief\\\emph{Machine Learning}\\Springer Nature}

\opening{Dear Editors,}

Please consider the enclosed manuscript, ``\emph{Beyond Planted Strings: Semantic
Canaries, Canary Detectability, and Partial Unlearning in Anytime-Valid Unlearning
Certificates},'' for publication in \emph{Machine Learning}.

\textbf{Relation to my earlier paper.} The manuscript extends ``\emph{VOUCH: Anytime-Valid
Certificates for Machine Unlearning via Paired-Canary Hypothesis-Betting},'' which has
been accepted by \emph{Machine Learning} [manuscript number: \emph{to be inserted}]. I
raise the overlap here so that you can judge it. The new manuscript recalls the VOUCH
certificate in one background section and cites the earlier paper for every result it
uses; it reproduces no theorem, table or figure of it. Everything stated as a result is
new: the theorems and propositions of its Sections~4--7, all experiments, and all
numbers. The two papers share the underlying construction (paired ghost canaries, the
exact null, betting e-processes), the fictitious-biography data generator, and the
released code base.

\textbf{What is new.} The earlier paper left three questions open, two of which it
named itself: whether canaries made of random strings say anything about the data that
users actually delete, what a provider that treats canaries differently could gain, and
whether an audit that reads the trained surface form sees knowledge that survives in
paraphrase. The manuscript addresses them as follows.
\begin{itemize}
\item \emph{Natural-language canaries.} It proves that semantic twins, two incompatible
facts about one fictitious entity, keep the exact null for every score that compares the
two facts under the same prompt, including belief scores over paraphrased questions that
were never trained and their versions calibrated by the base model.
\item \emph{From canaries to organic records.} Under a bounded-influence assumption the
cohort certificate transfers to the canary generator's population and then to organic
records at a cost of twice the provider's ability to tell the two apart from their
contents. For one-sided leakage the dependence is sharp: the largest organic advantage
compatible with a certified cohort is a Neyman--Pearson power. Two measurements go with
it: every canary design built for VOUCH, and naive semantic twins, is detected with
advantage one, and a provider guided by an off-the-shelf detector leaves essentially all
organic records un-unlearned while the cohort advantage stays at the tolerance.
\item \emph{The cost of enlarging the class and of monitoring continuously.} Adding
scores costs nothing in error on the certificate side and a logarithm on the revocation
side; in the sign model no sequential test needs fewer than $\log(1/\alpha)/\mathrm{KL}$
samples, and an oracle bettor reaches that within a logarithmic term, which matches a
fixed-sample sign test at about $78\%$ power.
\item \emph{Experiments on Pythia-410M and Pythia-1B} with eight seeds per setting, and
a comparison with membership-inference audits on the same runs.
\end{itemize}

\textbf{What the experiments do and do not show.} I want to be direct about the limits,
because several results are negative or qualified, and the manuscript reports them as such.
\begin{itemize}
\item The squeezing effect, that unlearning removes the literal form faster than the
fact, appears in the expected direction in all three experiments, but its size
($+0.07$ to $+0.12$ in retention, like for like) is not significant with eight seeds.
The manuscript does not claim to have demonstrated it. An earlier analysis of mine
overstated the effect by comparing a raw score with a calibrated one and by treating
strength-by-seed cells as replicates; both are corrected, and the manuscript says so.
\item Organic-matched canaries, which a verifier would choose to make the certificate
transfer to organic records, are not detected by the detectors I built, but a record seen
once leaves no detectable trace in the model at the training budget used, so the audit is
valid and powerless there. The manuscript presents this as the price of the transfer, not
as a flaw to be hidden.
\item A membership-inference test of canaries against an organic holdout raised a false
alarm on every retrained model; against matched ghost twins it is valid but weaker than a
paired test on the same records.
\item The theory is conditional on the distinguishing advantage, which data can bound
only from below, and on a bounded-influence assumption that cannot be checked. All facts
are fictitious, from one domain, and the models have at most one billion parameters.
\end{itemize}

\textbf{Reproducibility.} All code, every per-seed result and the scripts that turn the
results into the tables and the figure of the manuscript are released in the repository
below, in the directory \texttt{extension/}, together with a test suite for the
generators and the numerical checks of the new propositions.
\begin{center}
\href{https://github.com/vinhqdang/VOUCH-Verifiable-Online-Unlearning-Certification-via-Hypothesis-betting}{\texttt{github.com/vinhqdang/VOUCH-Verifiable-Online-}\\\texttt{Unlearning-Certification-via-Hypothesis-betting}}
\end{center}

This manuscript is original, has not been published previously, and is not under
consideration elsewhere. I am the sole author. I have no conflicts of interest to declare.
I would be glad to suggest reviewers, or to supply the earlier paper's referee reports and
my responses to them, if that would help you assess the overlap.

\closing{Sincerely,}

\end{letter}
\end{document}
